%==============================================================================
%  Interrogation n°2 - Mathematiques - Classe de 5e
%  Theme : Angles (adjacents, complementaires, supplementaires)
%  Compilation : pdflatex
%==============================================================================
\documentclass[12pt,a4paper]{article}

\usepackage[utf8]{inputenc}
\usepackage[T1]{fontenc}
\usepackage{lmodern}
\usepackage[french]{babel}
\usepackage[a4paper,margin=2cm]{geometry}
\usepackage{amsmath,amssymb}
\usepackage{array}
\usepackage{tabularx}            % colonne qui s'adapte a la largeur
\usepackage{booktabs}
\renewcommand{\arraystretch}{1.5}

% Colonne centree de largeur fixe
\newcolumntype{C}[1]{>{\centering\arraybackslash}p{#1}}

\pagestyle{empty}
\setlength{\parindent}{0pt}

\begin{document}

%------------------------------------------------------------------ EN-TETE
\begin{center}
\begin{tabularx}{\textwidth}{|X|C{4.2cm}|}
\hline
\textbf{Lycée Mixte 1 de Yamoussoukro} & \textbf{Année scolaire : 2025-2026} \\
\textbf{Discipline : Mathématiques} & \textbf{Classe : 5\textsuperscript{e}} \\
\textbf{Professeur : \dotfill} & \textbf{Durée : 15 min} \\
\hline
\end{tabularx}
\end{center}

\vspace{0.4cm}
\begin{center}
{\Large\bfseries INTERROGATION ÉCRITE N\textsuperscript{o} 2}\\[2pt]
{\large \textit{Thème : Les angles}}
\end{center}

\vspace{0.2cm}
\noindent\textbf{Nom et prénoms :} \dotfill \hfill \textbf{Note :} \fbox{\rule{1.2cm}{0pt}\rule{0pt}{0.6cm}} / 10

\vspace{0.5cm}
\hrule
\vspace{0.5cm}

%------------------------------------------------------------------ CONSIGNE
\textbf{Consigne.} Pour chacune des affirmations suivantes, mets une croix (\(\times\)) dans la case
\textbf{VRAI} si l'affirmation est vraie, ou dans la case \textbf{FAUX} si elle est fausse.
\textit{(Chaque réponse juste vaut 1 point.)}

\vspace{0.4cm}

%------------------------------------------------------------------ TABLEAU VRAI/FAUX
\begin{center}
\begin{tabularx}{\textwidth}{|C{0.8cm}|X|C{1.8cm}|C{1.8cm}|}
\hline
\textbf{N\textsuperscript{o}} & \centering\textbf{Affirmations} & \textbf{VRAI} & \textbf{FAUX} \\
\hline
1 & Deux angles adjacents ont le même sommet et un côté commun. & & \\
\hline
2 & Si deux angles ont seulement un sommet commun, alors ils sont adjacents. & & \\
\hline
3 & Deux angles adjacents sont situés de part et d'autre de leur côté commun. & & \\
\hline
4 & Deux angles complémentaires sont deux angles dont la somme des mesures est égale à 90\degres. & & \\
\hline
5 & Deux angles supplémentaires sont deux angles dont la somme des mesures est égale à 90\degres. & & \\
\hline
6 & Un angle de mesure 35\degres{} et un angle de mesure 55\degres{} sont complémentaires. & & \\
\hline
7 & Un angle de mesure 120\degres{} et un angle de mesure 70\degres{} sont supplémentaires. & & \\
\hline
8 & Le supplément d'un angle de mesure 60\degres{} a pour mesure 120\degres. & & \\
\hline
9 & Deux angles complémentaires peuvent être tous les deux obtus. & & \\
\hline
10 & Le complément d'un angle de mesure 45\degres{} a pour mesure 45\degres. & & \\
\hline
\end{tabularx}
\end{center}

\vfill
\begin{center}
\textit{« Le travail bien fait est la clé de la réussite. » Bonne chance !}
\end{center}

\end{document}
